% !TEX TS-program = lualatex
% !TEX encoding = UTF-8

\documentclass[antiphonaire_ordinaire.tex]{subfiles}

\ifcsname preamble@file\endcsname
  \setcounter{page}{\getpagerefnumber{M-ao_ordinarium}}
\fi

\begin{document}

\titulum{1}{Ordinaire de l'Office}{}{}

\titulum{2}{Aux Laudes}{Ordinaire}{Laudes}

\rubrica{Le dimanche et aux fêtes de troisième classe:}

\smallscore{DIA_festivus}{\vvtrans Dieu, viens à mon aide. \rrtrans Seigneur, viens vite à mon secours. Gloire au Père, et au Fils, et au Saint-Esprit, comme il était au commencement, maintenant et toujours, dans les siècles des siècles. Amen. Alléluia. \textcolor{gregoriocolor}{De la Septuagésime à la Semaine Sainte, on remplace \normaltext{Alléluia} par:} Louange à toi, ô Christ, roi d'éternelle gloire.}

\rubrica{Aux féries:}

\smallscore{DIA_ferialis}{}

\titulumrubrica{Psalmodie}

\rubrica{Antiennes et Psaumes, au Psautier.}

\titulumrubrica{Capitule, Hymne et Verset}

\rubrica{Le dimanche et aux féries, Capitule, Hymne et Verset du jour, au Psautier. Aux fêtes, au Propre des Saints, ou, à défaut, au Commun.}

\titulumrubrica{Cantique}

\rubrica{Le dimanche, Antienne à Benedictus au Propre du Temps. Aux féries, Antienne au Psautier. Aux fêtes, au Propre des Saints, ou, à défaut, au Commun.}

\renewcommand{\nextpsalmmode}{7a} % This is so the 2 accents pointing shows up next.
\canticum{Benedictus}{Cantique de Zacharie}{Lc. 1: 68-79}

\titulumrubrica{Conclusion}

\rubrica{Après avoir répété l'antienne du Benedictus, si le célébrant est au moins diacre:}

\smallscore{ORDV}{\vvtrans Le Seigneur soit avec vous. \rrtrans Et avec votre esprit. \vvtrans Prions.}

\rubrica{Si le célébrant n'est pas au moins diacre:}

\smallscore{ORDE}{\vvtrans Seigneur, entends ma prière. \rrtrans Que mon cri parvienne jusqu'à toi. \vvtrans Prions.}

\rubrica{Oraison au Propre du Temps ou au Propre des Saints. On répond \normaltext{Amen}, puis, sur le même ton que ci-dessus:}

\dominusvobiscumversiculus
\benedicamusversiculus
\fideliumversiculus

\rubrica{Tons du \normaltext{Benedicámus}:}

\rubrica{Le dimanche:}
\smallscore{ORBDF1}{}

\rubrica{Aux fêtes de 3\textsuperscript{e} classe:}
\smallscore{ORBDC3L}{}

\rubrica{Ou bien:}
\smallscore{ORBDC3X}{}

\rubrica{Aux féries:}
\smallscore{ORBDF2}{}{}

\titulum{2}{À Prime}{Ordinaire}{Prime}

\smallscore{DIA_ferialis}{\vvtrans Dieu, viens à mon aide. \rrtrans Seigneur, viens vite à mon secours. Gloire au Père, et au Fils, et au Saint-Esprit, comme il était au commencement, maintenant et toujours, dans les siècles des siècles. Amen. Alléluia. \textcolor{gregoriocolor}{De la Septuagésime à la Semaine Sainte, on remplace \normaltext{Alléluia} par:} Louange à toi, ô Christ, roi d'éternelle gloire.}

\titulumrubrica{Hymne}

\rubrica{Le dimanche et aux fêtes de troisième classe:}
\hymnus{ORPHb}

\rubrica{Aux féries:}
\gscore{ORPHc}{}{}

\titulumrubrica{Psalmodie}

\rubrica{Antienne et Psaumes, au Psautier.}

\titulumrubrica{Capitule, Répons bref et Verset}
\label{ORP}

\capitulum{1 Tim. 1: 17}{Regi sæculórum immortáli et invisíbili,\pscross{} soli Deo honor et glória\psstar{} in sǽcula sæculórum. Amen. ~~~~\rr Deo grátias.}{Au roi des siècles, Dieu immortel, invisible et unique, honneur et gloire pour les siècles des siècles ! Amen. ~~~~\rr Nous rendons grâces à Dieu.}

\gscore{ORPRa}{}{\rrtrans Christ, Fils du Dieu vivant, aie pitié de nous. \vvtrans Toi qui es assis à la droite du Père. \vvtrans Gloire au Père, au Fils, et au Saint-Esprit.}

\rubrica{De Noël au 5 janvier, en la fête du Corps du Christ, aux fêtes de la Sainte Vierge:}
\smallscore{ORPR_quinatuses}{\vvtrans Toi qui es né de la Vierge Marie.}

\rubrica{De l'Épiphanie au 13 janvier:}
\smallscore{ORPR_quiapparuisti}{\vvtrans Toi qui t'es manifesté aujourd'hui.}

\versiculus{Exsúrge, Christe, ádjuva nos.}{Et líbera nos propter nomen tuum.}{Lève-toi, ô Christ, aide-nous.}{Et délivre-nous pour la gloire de ton nom.}


\dominusvobiscumversiculus

\rubrica{Si le célébrant n'est pas au moins diacre:}

\dominexaudiversiculus

\versiculus{Orémus.\\Dómine Deus omnípotens, qui ad princípium hujus diéi nos perveníre fe\sylac{cí}sti:\pscross{} tua nos hódie salva virtúte; ut in hac die ad nullum declinémus peccátum, sed semper ad tuam justítiam faciéndam nostra procé\sylac{dant} \sylac{e}\sylprep{ló}quia,\psstar{} dirigántur cogitatiónes et ópera. Per Dóminum nostrum Jesum Christum, Fílium \sylprep{tu}um:\pscross{} qui tecum vivit et regnat in unitáte Spíritus \sylac{Sancti} \sylprep{De}us,\psstar{} per ómnia sǽcula sæculórum.}{Amen.}{Prions.\\Seigneur, Dieu tout-puissant, qui nous as fait parvenir au commencement de ce jour, sauve-nous aujourd'hui par ta puissance, afin que durant ce jour, nous ne nous cédions à aucun péché : mais que nos paroles, nos pensées et nos œuvres tendent toujours à l'accomplissement de ta justice. Par Jésus-Christ, ton Fils, notre Seigneur, qui vit et règne avec toi dans l'unité du Saint-Esprit, Dieu, pour les siècles des siècles.}{Amen.}

\dominusvobiscumversiculus

\rubrica{Ou bien \normaltext{Dómine, exáudi}, etc.}

\smallscore{ORBDhm}{\vvtrans Bénissons le Seigneur. \rrtrans Nous rendons grâces à Dieu.}

\titulumrubrica{À l'office capitulaire}

\rubrica{Le lecteur chante le Martyrologe du jour, et termine par:}
\versiculus{Et álibi aliórum plurimórum sanctórum Mártyrum et Confessórum, atque sanctárum Vírginum.}{Deo grátias.}{Et ailleurs, beaucoup d’autres saints martyrs, confesseurs et saintes vierges.}{Nous rendons grâces à Dieu.}

\versiculus{Pretiósa in conspéctu Dómini.}{Mors Sanctórum ejus.}{Elle est précieuse aux yeux du Seigneur.}{La mort de ses saints.}
\versiculus{Sancta María et omnes Sancti intercédant pro nobis ad \sylac{Dó}minum,\pscross{} ut nos mereámur ab eo adjuvári \sylprep{et} \sylprep{sal}vári,\psstar{} qui vivit et regnat in sǽcula sæculórum.}{Amen.}{Que sainte Marie et tous les saints intercèdent pour nous auprès du Seigneur, afin qu'il nous secoure et nous sauve, lui qui vit et règne pour les siècles des siècles.}{Amen.}

\versiculus{Deus in adjutórium meum inténde.}{Dómine, ad adjuvándum me festína.}{Dieu, viens à mon aide.}{Seigneur, viens vite à mon secours}
\versiculus{Deus in adjutórium meum inténde.}{Dómine, ad adjuvándum me festína.}{Dieu, viens à mon aide.}{Seigneur, viens vite à mon secours}
\versiculus{Deus in adjutórium meum inténde.}{Dómine, ad adjuvándum me festína.}{Dieu, viens à mon aide.}{Seigneur, viens vite à mon secours}
\versiculus{Glória Pa\sylprep{tri,} \sylprep{et} \sylac{Fí}lio,\psstar{} et Spirítui Sancto.}{Sicut erat in princípio, et \sylprep{nunc,} \sylprep{et} \sylac{sem}per,\psstar{} et in sǽcula sæculórum. Amen.}{Gloire au Père, et au Fils, et au Saint-Esprit.}{Comme il était au commencement, maintenant et toujours, dans les siècles des siècles. Amen.}

\smallscore{ORPK}{}

\rubrica{en silence jusqu'à \normaltext{Et ne nos}.}

\twocoltext{Pater noster, qui es in cælis, sanctificétur nomen tuum: advéniat regnum tuum: fiat volúntas tua, sicut in cælo et in terra. Panem nostrum quotidiánum da nobis hódie: et dimítte nobis débita nostra, sicut et nos dimíttimus debitóribus nostris:}{Notre Père, qui es au cieux, que ton nom soit sanctifié; que ton règne vienne; que ta volonté soit faite sur la terre comme au ciel. Donne-nous aujourd'hui notre pain de ce jour; pardonne-nous nos offenses, comme nous pardonnons aussi à ceux qui nous ont offensés.}

\smallscore{ORPP2}{\vvtrans Et ne nous laisse pas entrer en tentation. \rrtrans Mais délivre-nous du mal.}

\versiculus{Réspice in servos tuos, Dómine, et in ó\sylprep{pera} \sylac{tu}a,\psstar{} et dírige fílios eórum.}{Et sit splendor Dómini Dei nostri \sylac{su}per nos,\pscross{} et ópera mánuum nostrárum dí\sylprep{rige} \sylac{su}per nos,\psstar{} et opus mánuum nostrárum dírige.}{Fais connaître ton œuvre à tes serviteurs, Seigneur, et ta splendeur à leurs fils.}{Que vienne sur nous la douceur du Seigneur notre Dieu! Consolide pour nous l'ouvrage de nos mains; oui, consolide l'ouvrage de nos mains.}
\versiculus{Glória Pa\sylprep{tri,} \sylprep{et} \sylac{Fí}lio,\psstar{} et Spirítui Sancto.}{Sicut erat in princípio, et \sylprep{nunc,} \sylprep{et} \sylac{sem}per,\psstar{} et in sǽcula sæculórum. Amen.}{Gloire au Père, et au Fils, et au Saint-Esprit.}{Comme il était au commencement, maintenant et toujours, dans les siècles des siècles. Amen.}
\versiculus{Orémus.\\
Dirígere et sanctificáre, régere et gubernáre dignáre, Dómine Deus, Rex cæli et terræ, hódie corda et córpora \sylac{nos}tra,\pscross{} sensus, sermónes et actus nostros in lege tua, et in opéribus mandató\sylprep{rum} \sylprep{tu}\sylac{ó}rum:\psstar{} ut hic et in ætérnum, te auxiliánte, salvi et líberi esse mereámur, Salvátor mundi: Qui vivis et regnas in sǽcula sæculórum.}{Amen.}{Prions.\\
Seigneur Dieu, Roi du ciel et de la terre, dirige et sanctifie, régis et gouverne en ce jour nos cœurs et nos corps, nos sens, nos paroles et nos actes, selon ta loi et l'accomplissement de tes préceptes, afin qu'ici-bas et pour l'éternité, nous obtenions d'être sauvés et délivrés par ton secours, ô Sauveur du monde: toi qui vis et règnes pour les siècles des siècles.}{Amen.}

\smallscore{ORPBenedictio}{Veuillez, père, bénir. \textcolor{gregoriocolor}{Bénédiction.} Que le Seigneur tout-puissant établisse dans Sa paix nos jours et nos actions.}
\rubrica{Si le célébrant n'est pas diacre, le lecteur s'incline vers la croix et non vers lui, et dit \normaltext{Jube, Dómine, benedícere}.}

%\rubrica{Pendant l'année, et au temps de la Septuagésime:}

\lbrev{2 Th. 3: 5}{Dóminus autem dírigat corda et córpora nostra in cari\sylprep{táte} \sylac{De}i\psstar{} et patiéntia Christi.\\Tu autem, Dómine, miserére nobis. ~~~~\rr Deo grátias.}{Que le Seigneur conduise nos cœurs dans l’amour de Dieu et l’endurance du Christ.\\Et toi, Seigneur, aie pitié de nous. ~~~~\rr Nous rendons grâces à Dieu.}

%\smallscore{ORTuAutem}{\vvtrans Et toi Seigneur, prends pitié de nous.}{\rrtrans Nous rendons grâces à Dieu.}

%\rubrica{Pendant l'Avent:}

%\lbrev{Is. 33: 2}{Dómine, miserére nostri: te enim exspec\sylac{tá}vimus:\pscross{} esto brácchium nos\sylprep{trum} \sylprep{in} \sylac{ma}ne,\psstar{} et salus nostra in témpore tribulatiónis.}{Seigneur, fais-nous grâce : c’est toi que nous attendons ! Chaque matin, sois notre bras, notre salut aux jours de détresse.}

%\rubrica{En Carême:}

%\lbrev{Is. 55: 6}{Quǽrite Dóminum dum inve\sylprep{níri} \sylac{pot}est:\psstar{} invocáte eum, dum prope est.}{Cherchez le Seigneur tant qu’il se laisse trouver ; invoquez-le tant qu’il est proche.}

%\rubrica{Au temps de la Passion:}

%\lbrev{Is. 50: 6-7}{Fáciem meam non avérti ab increpántibus, et conspuéntibus \sylac{in} me.\pscross{} Dóminus Deus auxili\sylprep{átor} \sylac{me}us,\psstar{} et ídeo non sum confúsus.}{Je n’ai pas caché ma face devant les outrages et les crachats. Le Seigneur mon Dieu vient à mon secours ; c’est pourquoi je ne suis pas atteint par les outrages.}

\versiculus{Adjutórium nostrum \cc in nómine Dómini.}{Qui fecit cælum et terram.}{Notre secours est le nom du Seigneur.}{Qui a fait le ciel et la terre.}
\versiculus{Benedícite.}{Deus.}{Bénissez.}{Ô Dieu.}
\smallscore{ORPBenFinal}{Que le Seigneur tout-puissant nous bénisse, nous protège de tout mal et nous conduise à la vie éternelle. Et que par la miséricorde de Dieu, les âmes des fidèles trépassés reposent en paix. \rrtrans Amen.}

\titulum{2}{À Tierce}{Ordinaire}{Tierce}

\smallscore{DIA_ferialis}{\vvtrans Dieu, viens à mon aide. \rrtrans Seigneur, viens vite à mon secours. Gloire au Père, et au Fils, et au Saint-Esprit, comme il était au commencement, maintenant et toujours, dans les siècles des siècles. Amen. Alléluia. \textcolor{gregoriocolor}{De la Septuagésime à la Semaine Sainte, on remplace \normaltext{Alléluia} par:} Louange à toi, ô Christ, roi d'éternelle gloire.}

\titulumrubrica{Hymne}

\rubrica{Le dimanche et aux fêtes de troisième classe:}
\hymnus{ORTHb}

\rubrica{Aux féries:}
\gscore{ORTHc}{}{}

\titulumrubrica{Psalmodie}

\rubrica{Antienne et Psaumes, au Psautier.}

\titulumrubrica{Capitule, Répons bref et Verset}
\label{ORT}

\rubrica{Le dimanche:}

\capitulum{1 Jn. 4: 16}{Deus cáritas est: et qui manet in caritáte, in Deo manet, et Deus in eo.}{Dieu est amour : qui demeure dans l’amour demeure en Dieu, et Dieu demeure en lui.}

\gscore{ORTRa}{}{\rrtrans Incline mon cœur, ô Dieu, * vers tes témoignages. \vvtrans Détourne mes yeux de la vue de la vanité ; dans ta voie, donne-moi la vie.}

\versiculus{Ego dixi: Dómine, miserére mei.}{Sana ánimam meam, quia peccávi tibi.}{J’ai dit : Seigneur, prends pitié de moi.}{Guéris mon âme parce que j’ai péché contre toi.}

\rubrica{Aux féries:}

\capitulum{Jr. 17: 14}{Sana me, Dómine, et sa\sylac{ná}bor:\pscross{} salvum me fac, et \sylprep{salvus} \sylac{e}ro:\psstar{} quóniam laus mea tu es.}{Guéris-moi, Seigneur, et je serai guéri, sauve-moi, et je serai sauvé, car tu es ma louange.}

\gscore{ORTRb}{}{\rrtrans Guéris mon âme, car j'ai péché contre toi. \vvtrans J'ai dit: Seigneur, prends pitié de moi.}

\versiculus{Adjútor meus, esto, ne derelínquas me.}{Neque despícias me, Deus, salutáris meus.}{Sois mon assistance, ne m'abandonne pas.}{Ne me méprise pas, Dieu de mon salut.}

\rubrica{Aux fêtes, Capitule, Répons bref et Verset au Propre ou au Commun.}

\titulumrubrica{Conclusion}

\dominusvobiscumversiculus

\rubrica{ou bien, si le célébrant n'est pas au moins diacre:}

\dominexaudiversiculus

\rubrica{Oraison du jour, ou, aux féries, du dimanche précédent.}

\dominusvobiscumversiculus

\rubrica{ou bien \normaltext{Dómine, exáudi}, etc.}

\smallscore{ORBDhm}{\vvtrans Bénissons le Seigneur. \rrtrans Nous rendons grâces à Dieu.}

\fideliumversiculus

\titulum{2}{À Sexte}{Ordinaire}{Sexte}

\titulumrubrica{Hymne}

\rubrica{Le dimanche et aux fêtes de troisième classe:}
\hymnus{ORSHb}

\rubrica{Aux féries:}
\gscore{ORSHc}{}{}

\titulumrubrica{Psalmodie}

\rubrica{Antienne et Psaumes, au Psautier.}

\titulumrubrica{Capitule, Répons bref et Verset}
\label{ORS}

\rubrica{Le dimanche:}

\capitulum{Ga. 6: 2}{Alter altérius óne\sylprep{ra} \sylprep{por}\sylac{tá}te,\psstar{} et sic adimplébitis legem Christi.}{Portez les fardeaux les uns des autres : ainsi vous accomplirez la loi du Christ.}

\gscore{ORSRa}{}{\rrtrans À jamais, Seigneur, demeure ta parole. \vvtrans Dans les siècles des siècles, ta vérité.}

\versiculus{Dóminus regit me, et nihil mihi déerit.}{In loco páscuæ ibi me collocávit.}{Le Seigneur est mon berger, je ne manque de rien.}{Sur des prés d'herbe fraîche, il me fait reposer.}

\rubrica{Aux féries:}

\capitulum{Rm. 13: 8}{Némini quidquam debe\sylac{á}tis,\pscross{} nisi ut ínvicem \sylprep{dili}\sylac{gá}tis:\psstar{} qui enim díligit próximum, legem implévit.}{N’ayez de dette envers personne, sauf celle de l’amour mutuel, car celui qui aime les autres a pleinement accompli la Loi.}

\gscore{ORSRb}{}{\rrtrans Je bénirai le Seigneur en tout temps \vvtrans Sa louange sans cesse à ma bouche.}

\versiculus{Dóminus regit me, et nihil mihi déerit.}{In loco páscuæ ibi me collocávit.}{Le Seigneur est mon berger, je ne manque de rien.}{Sur des prés d'herbe fraîche, il me fait reposer.}

\rubrica{Aux fêtes, Capitule, Répons bref et Verset au Propre ou au Commun.}

\titulumrubrica{Conclusion}

\dominusvobiscumversiculus

\rubrica{ou bien, si le célébrant n'est pas au moins diacre:}

\dominexaudiversiculus

\rubrica{Oraison du jour, ou, aux féries, du dimanche précédent.}

\dominusvobiscumversiculus

\rubrica{ou bien \normaltext{Dómine, exáudi}, etc.}

\smallscore{ORBDhm}{\vvtrans Bénissons le Seigneur. \rrtrans Nous rendons grâces à Dieu.}

\fideliumversiculus

\titulum{2}{À None}{Ordinaire}{None}

\smallscore{DIA_ferialis}{\vvtrans Dieu, viens à mon aide. \rrtrans Seigneur, viens vite à mon secours. Gloire au Père, et au Fils, et au Saint-Esprit, comme il était au commencement, maintenant et toujours, dans les siècles des siècles. Amen. Alléluia. \textcolor{gregoriocolor}{De la Septuagésime à la Semaine Sainte, on remplace \normaltext{Alléluia} par:} Louange à toi, ô Christ, roi d'éternelle gloire.}

\titulumrubrica{Hymne}

\rubrica{Le dimanche et aux fêtes de troisième classe:}
\hymnus{ORNHb}

\rubrica{Aux féries:}
\gscore{ORNHc}{}{}

%% TODO ajouter les tons de l'hymne du temps

\titulumrubrica{Psalmodie}

\rubrica{Antienne et Psaumes, au Psautier.}

\titulumrubrica{Capitule, Répons bref et Verset}
\label{ORN}

\rubrica{Le dimanche:}

\capitulum{1 Co. 6: 20}{Empti enim estis pré\sylprep{tio} \sylac{ma}gno.\psstar{} Glorificáte et portáte Deum in córpore vestro.}{Vous avez été achetés à grand prix. Rendez donc gloire à Dieu dans votre corps.}

\gscore{ORNRa}{}{\rrtrans J’ai crié de tout mon cœur: exauce-moi, Seigneur. \vvtrans C'est ta justice que je cherche.}

\versiculus{Ab occúltis meis munda me, Dómine.}{Et ab aliénis parce servo tuo.}{Purifie-moi de mes fautes cachées, Seigneur.}{Et épargne-moi une domination étrangère.}

\rubrica{Aux féries:}

\capitulum{1 P. 1: 17-19}{In timóre incolátus vestri témpore conver\sylac{sá}mini:\pscross{} sciéntes quod non corruptibílibus auro vel argénto red\sylprep{émpti} \sylac{es}tis,\psstar{} sed pretióso sánguine quasi Agni immaculáti Christi.
}{Vivez dans la crainte de Dieu, pendant le temps où vous résidez ici-bas en étrangers. Vous le savez : ce n’est pas par des biens corruptibles, l’argent ou l’or, que vous avez été rachetés de la conduite superficielle héritée de vos pères ; mais c’est par un sang précieux, celui d’un agneau sans défaut, le Christ.}

\gscore{ORNRb}{}{\rrtrans Délivre-moi, Seigneur, et prends pitié de moi. \vvtrans Mon pied s'est tenu dans la voie droite.}

\versiculus{Ab occúltis meis munda me, Dómine.}{Et ab aliénis parce servo tuo.}{Purifie-moi de mes fautes cachées, Seigneur.}{Et épargne-moi une domination étrangère.}

\rubrica{Aux fêtes, Capitule, Répons bref et Verset au Propre ou au Commun.}

\titulumrubrica{Conclusion}

\dominusvobiscumversiculus

\rubrica{ou bien, si le célébrant n'est pas au moins diacre:}

\dominexaudiversiculus

\rubrica{Oraison du jour, ou, aux féries, du dimanche précédent.}

\dominusvobiscumversiculus

\rubrica{ou bien \normaltext{Dómine, exáudi}, etc.}

\smallscore{ORBDhm}{\vvtrans Bénissons le Seigneur. \rrtrans Nous rendons grâces à Dieu.}

\fideliumversiculus

\titulum{2}{Aux Vêpres}{Ordinaire}{Vêpres}

\rubrica{Le dimanche et aux fêtes de troisième classe:}

\smallscore{DIA_festivus}{\vvtrans Dieu, viens à mon aide. \rrtrans Seigneur, viens vite à mon secours. Gloire au Père, et au Fils, et au Saint-Esprit, comme il était au commencement, maintenant et toujours, dans les siècles des siècles. Amen. Alléluia. \textcolor{gregoriocolor}{De la Septuagésime à la Semaine Sainte, on remplace \normaltext{Alléluia} par:} Louange à toi, ô Christ, roi d'éternelle gloire.}

\rubrica{Aux féries:}

\smallscore{DIA_ferialis}{}

\titulumrubrica{Psalmodie}

\rubrica{Antiennes et Psaumes, au Psautier.}

\titulumrubrica{Capitule, Hymne et Verset}

\rubrica{Le dimanche et aux féries, Capitule, Hymne et Verset du jour, au Psautier. Aux fêtes, au Propre des Saints, ou, à défaut, au Commun.}

\titulumrubrica{Cantique}

\rubrica{Le dimanche, Antienne à Magnificat au Propre du Temps. Aux féries, Antienne au Psautier. Aux fêtes, au Propre des Saints, ou, à défaut, au Commun.}

\renewcommand{\nextpsalmmode}{7a} % This is so the 2 accents pointing shows up next.
\canticum{Magnificat}{Cantique de la Bienheureuse Vierge Marie}{Lc. 1: 46-55}

\titulumrubrica{Conclusion}

\rubrica{Après avoir répété l'antienne du Magnificat, si le célébrant est au moins diacre:}

\smallscore{ORDV}{\vvtrans Le Seigneur soit avec vous. \rrtrans Et avec votre esprit. \vvtrans Prions.}

\rubrica{Si le célébrant n'est pas au moins diacre:}

\smallscore{ORDE}{\vvtrans Seigneur, entends ma prière. \rrtrans Que mon cri parvienne jusqu'à toi. \vvtrans Prions.}

\rubrica{Oraison au Propre du Temps ou au Propre des Saints. On répond \normaltext{Amen}, puis, sur le même ton que ci-dessus:}

\dominusvobiscumversiculus
\benedicamusversiculus
\fideliumversiculus

\rubrica{Tons du \normaltext{Benedicámus}:}

\rubrica{Le dimanche:}
\smallscore{ORBDF1}{}

\rubrica{Aux fêtes de 3\textsuperscript{e} classe:}
\smallscore{ORBDC3V}{}

\rubrica{Ou bien:}
\smallscore{ORBDC3X}{}

\rubrica{Aux féries:}
\smallscore{ORBDF2}{}{}

\titulum{2}{Aux Complies}{Ordinaire}{Complies}

\smallscore{ORCBenedictio}{Veuillez, père, bénir. \textcolor{gregoriocolor}{Bénédiction.} Que le Seigneur tout-puissant nous accorde une nuit tranquille et une fin parfaite.}
\rubrica{Si le célébrant n'est pas diacre, le lecteur s'incline vers la croix et non vers lui, et dit \normaltext{Jube, Dómine, benedícere}.}

\lbrev{1 P. 5: 8-9}{Fratres: Sóbrii estóte, et vigi\sylac{lá}te:\pscross{} quia adversárius vester diábolus tamquam leo rúgiens círcuit, quæ\sylprep{rens} \sylprep{quem} \sylac{dé}voret:\psstar{} cui resístite fortes in fide.\\Tu autem, Dómine, miserére nobis. ~~~~\rr Deo grátias.}{Frères, soyez sobres, veillez: votre adversaire, le diable, comme un lion rugissant, rôde, cherchant qui dévorer. Résistez-lui avec la force de la foi.\\Et toi, Seigneur, aie pitié de nous. ~~~~\rr Nous rendons grâces à Dieu.}

\smallscore{ORC_AdjutoriumNostrum}{\vvtrans Notre secours est le nom du Seigneur. \rrtrans Qui a fait le ciel et la terre.}
\rubrica{Examen de conscience en silence, le temps d'un \normaltext{Pater}.}

\rubrica{Si le célébrant est au moins diacre, il récite:}

\twocoltext{Confíteor Deo omnipoténti, beátæ Maríæ semper Vírgini, beáto Michaéli Archángelo, beáto Joánni Baptístæ, sanctis Apóstolis Petro et Paulo, ómnibus Sanctis, et vobis fratres, quia peccávi nimis, cogitatióne, verbo et ópere:}{Je confesse à Dieu tout puissant, à la bienheureuse Marie toujours Vierge, à saint Michel Archange, à saint Jean Baptiste, aux saints Apôtres Pierre et Paul, à tous les saints, et à vous, mes frères, que j'ai beaucoup péché, en pensées, en paroles, et par actions.}
\rubrica{Il se frappe la poitrine.}
\twocoltext{Mea culpa, mea culpa, mea máxima culpa. Ídeo precor beátam Maríam semper Vírginem, beátum Michaélem Archángelum, beátum Joánnem Baptístam, sanctos Apóstolos Petrum et Paulum, omnes Sanctos, et vos fratres, oráre pro me ad Dóminum Deum nostrum.}{C'est ma faute, c'est ma faute, c'est ma très grande faute. C'est pourquoi je supplie la bienheureuse Marie toujours Vierge, saint Michel Archange, saint Jean Baptiste, les saints Apôtres Pierre et Paul, tous les saints, et vous, mes frères, de prier pour moi le Seigneur notre Dieu.}
\rubrica{Tous répondent:}
\twocoltext{Misereátur tui omnípotens Deus, et dimíssis peccátis nostris, perdúcat nos ad vitam ætérnam. \rr Amen.}{Que Dieu tout-puissant nous fasse miséricorde, qu'il nous pardonne nos péchés et nous conduise à la vie éternelle. \rr Amen.}
\rubrica{Tous récitent:}
\twocoltext{Confíteor Deo omnipoténti, beátæ Maríæ semper Vírgini, beáto Michaéli Archángelo, beáto Joánni Baptístæ, sanctis Apóstolis Petro et Paulo, ómnibus Sanctis, et tibi pater, quia peccávi nimis, cogitatióne, verbo et ópere:}{Je confesse à Dieu tout puissant, à la bienheureuse Marie toujours Vierge, à saint Michel Archange, à saint Jean Baptiste, aux saints Apôtres Pierre et Paul, à tous les saints, et à vous, mon père, que j'ai beaucoup péché, en pensées, en paroles, et par actions.}
\rubrica{On se frappe la poitrine.}
\twocoltext{Mea culpa, mea culpa, mea máxima culpa. Ídeo precor beátam Maríam semper Vírginem, beátum Michaélem Archángelum, beátum Joánnem Baptístam, sanctos Apóstolos Petrum et Paulum, omnes Sanctos, et te pater, oráre pro me ad Dóminum Deum nostrum.}{C'est ma faute, c'est ma faute, c'est ma très grande faute. C'est pourquoi je supplie la bienheureuse Marie toujours Vierge, saint Michel Archange, saint Jean Baptiste, les saints Apôtres Pierre et Paul, tous les saints, et vous, mon père, de prier pour moi le Seigneur notre Dieu.}
\versiculus{Misereátur nostri omnípotens Deus, et dimíssis peccátis nostris, perdúcat nos ad vitam ætérnam.}{Amen.}{Que Dieu tout-puissant nous fasse miséricorde, qu'il nous pardonne nos péchés et nous conduise à la vie éternelle.}{Amen.}
\versiculus{Indulgéntiam, \cc absolutiónem et remissiónem peccatórum nostrórum tríbuat nobis omnípotens et miséricors Dóminus.}{Amen.}{Que le Seigneur tout-puissant et miséricordieux nos accorde l'indulgence, l'absolution et la rémission de nos péchés.}{Amen.}

\rubrica{Si le célébrant n'est pas au moins diacre, tous récitent ensemble:}

\twocoltext{Confíteor Deo omnipoténti, beátæ Maríæ semper Vírgini, beáto Michaéli Archángelo, beáto Joánni Baptístæ, sanctis Apóstolis Petro et Paulo, et ómnibus Sanctis, quia peccávi nimis, cogitatióne, verbo et ópere:}{Je confesse à Dieu tout puissant, à la bienheureuse Marie toujours Vierge, à saint Michel Archange, à saint Jean Baptiste, aux saints Apôtres Pierre et Paul, et à tous les saints, que j'ai beaucoup péché, en pensées, en paroles, et par actions.}
\rubrica{On se frappe la poitrine.}
\twocoltext{Mea culpa, mea culpa, mea máxima culpa. Ídeo precor beátam Maríam semper Vírginem, beátum Michaélem Archángelum, beátum Joánnem Baptístam, sanctos Apóstolos Petrum et Paulum, et omnes Sanctos, oráre pro me ad Dóminum Deum nostrum.}{C'est ma faute, c'est ma faute, c'est ma très grande faute. C'est pourquoi je supplie la bienheureuse Marie toujours Vierge, saint Michel Archange, saint Jean Baptiste, les saints Apôtres Pierre et Paul, et tous les saints, de prier pour moi le Seigneur notre Dieu.}
\versiculus{Misereátur nostri omnípotens Deus, et dimíssis peccátis nostris, perdúcat nos ad vitam ætérnam.}{Amen.}{Que Dieu tout-puissant nous fasse miséricorde, qu'il nous pardonne nos péchés et nous conduise à la vie éternelle.}{Amen.}
\versiculus{Indulgéntiam, \cc absolutiónem et remissiónem peccatórum nostrórum tríbuat nobis omnípotens et miséricors Dóminus.}{Amen.}{Que le Seigneur tout-puissant et miséricordieux nos accorde l'indulgence, l'absolution et la rémission de nos péchés.}{Amen.}

\rubrica{On fait un signe de croix sur sa poitrine.}

\smallscore{ORC_ConverteNos}{\vvtrans Convertis-nous, Dieu, notre Sauveur. \rrtrans Et détourne de nous ta colère.}

\smallscore{DIA_ferialis}{\vvtrans Dieu, viens à mon aide. \rrtrans Seigneur, viens vite à mon secours. Gloire au Père, et au Fils, et au Saint-Esprit, comme il était au commencement, maintenant et toujours, dans les siècles des siècles. Amen. Alléluia. \textcolor{gregoriocolor}{De la Septuagésime à la Semaine Sainte, on remplace \normaltext{Alléluia} par:} Louange à toi, ô Christ, roi d'éternelle gloire.}

\titulumrubrica{Psalmodie}

\rubrica{Antienne et Psaumes, au Psautier.}

\titulumrubrica{Hymne, Répons bref et Verset}

\rubrica{Le dimanche et aux fêtes de troisième classe:}

\hymnus{ORCHb}

\rubrica{Aux féries pendant l'année:}

\gscore{ORCHc}{}{}

\rubrica{Pendant le temps de Noël:}

\gscore{ORCHe}{}{}

\rubrica{Pendant le temps de l'Épiphanie:}

\gscore{ORCHf}{}{}

\gscore{ORCRa}{}{\rrtrans En tes mains, Seigneur, je remets mon esprit. \vvtrans Tu nous rachètes, Seigneur, Dieu de vérité.}

\versiculus{Custódi nos, Dómine, ut pupíllam óculi.}{Sub umbra alárum tuárum prótege nos.}{Garde-moi comme la prunelle de l'oeil.}{À l'ombre de tes ailes, cache-moi.}

%TODO: faut-il inclure les tons du In manus tuas?

\titulumrubrica{Cantique}

\gscore{ORCNAa}{}{Sauve-nous Seigneur, quand nous veillons, garde-nous quand nous dormons, nous veillerons avec le Christ et nous reposerons en paix.}
\canticum{NuncDimittis}{Cantique de Siméon}{Lc. 2: 29-32}

\titulumrubrica{Conclusion}

\dominusvobiscumversiculus

\rubrica{ou bien, si le célébrant n'est pas au moins diacre:}

\dominexaudiversiculus

\versiculus{Orémus.\\
Vísita, quǽsumus, Dómine, habitatiónem istam, et omnes insídias inimíci ab ea longe repélle: Ángeli tui sancti hábitent in ea, qui nos in pace custódiant; et benedíctio tua sit super nos semper. Per Dóminum nostrum Jesum Christum, Fílium tuum: qui tecum vivit et regnat in unitáte Spíritus Sancti, Deus, per ómnia sǽcula sæculórum.}{Amen.}{Prions.\\
Nous t’en supplions, Seigneur, visite cette maison, et repousse loin d’elle toutes les embûches de l’ennemi ; que tes saints anges viennent l’habiter pour nous garder dans la paix ; et que ta bénédiction demeure à jamais sur nous. Par Jésus Christ, ton Fils, notre Seigneur, qui vit et règne avec toi dans l'unité du Saint-Esprit, Dieu, pour les siècles des siècles.}{Amen.}

\dominusvobiscumversiculus

\rubrica{ou bien \normaltext{Dómine, exáudi}, etc.}

\smallscore{ORBDhm}{\vvtrans Bénissons le Seigneur. \rrtrans Nous rendons grâces à Dieu.}

\smallscore{ORCBenFinal}{Que le Seigneur tout-puissant et miséricordieux, \cc Père, Fils et Saint-Esprit, nous bénisse et nous garde. \rrtrans Amen.}

\titulumrubrica{Antienne finale à la Sainte Vierge}

\rubrica{Du 14 janvier au 1er février inclus:}

\gscore{ORC_AlmaA}{}{Sainte Mère du Rédempteur, porte du ciel, toujours ouverte, étoile de la mer, viens au secours du peuple qui tombe et qui cherche à se relever. Tu as enfanté, ô merveille ! celui qui t’a créée, et tu demeures toujours vierge. Accueille le salut de l’ange Gabriel et prends pitié de nous, pécheurs.}

\rubrica{ton simple:}
\gscore{ORC_AlmaB}{}{}

\versiculus{Post partum, Virgo, invioláta permansísti.}{Dei Génitrix, intercéde pro nobis.}{Tu es restée vierge après l'enfantement.}{Mère de Dieu, intercède pour nous.}
\versiculus{Orémus.\\
Deus, qui salútis ætérnæ, beátæ Maríæ virginitáte fecúnda, humáno géneri prǽmia præsti\sylac{tí}sti:\pscross{} tríbue, quǽsumus; ut ipsam pro nobis intercédere \sylprep{senti}\sylac{á}mus,\psstar{} per quam merúimus auctórem vitæ suscípere, Dóminum nostrum Jesum Christum Fílium tuum.}{Amen.}{Prions.\\
Seigneur Dieu, par la virginité féconde de la bienheureuse Marie, tu as offert au genre humain les bienfaits du salut éternel ; accorde-nous d’éprouver qu’intercède en notre faveur celle qui nous permit d’accueillir l’auteur de la vie, Jésus Christ, ton Fils, notre Seigneur.}{Amen.}

\versiculus{Divínum auxílium \cc máneat semper nobíscum.}{Amen.}{Que le secours divin demeure toujours avec nous.}{Amen.}

\rubrica{Du 2 février au Carême:}

\gscore{ORC_AveReginaA}{}{Salut, Reine des cieux ! Salut, Reine des anges ! Salut, Tige féconde ! Salut, Porte du ciel ! Par toi, la lumière s’est levée sur le monde. Réjouis-toi, Vierge glorieuse, belle entre toutes les femmes ! Salut, splendeur radieuse : implore le Christ pour nous.}

\rubrica{ton simple:}
\gscore{ORC_AveReginaB}{}{}

\versiculus{Dignáre me laudáre te, Virgo sacráta.}{Da mihi virtútem contra hostes tuos.}{Rends-moi digne de te louer, Vierge sainte.}{Donne-moi de la force contre tes ennemis.}
\versiculus{Orémus.\\
Concéde, miséricors Deus, fragilitáti nostræ præ\sylac{sí}dium;\pscross{} ut, qui sanctæ Dei Genitrícis memó\sylprep{riam} \sylac{á}gimus, intercessiónis ejus auxílio, a nostris iniquitátibus resurgámus. Per eúmdem Christum Dóminum nostrum.}{Amen.}{Prions.\\
Dieu de miséricorde, accorde ton secours à notre fragilité, afin que, en faisant mémoire de la sainte Mère de Dieu, et munis de son intercession, nous nous relèvions de nos péchés. Par le même Jésus, le Christ, notre Seigneur.}{Amen.}

\versiculus{Divínum auxílium \cc máneat semper nobíscum.}{Amen.}{Que le secours divin demeure toujours avec nous.}{Amen.}

\rubrica{Au temps après la Pentecôte:}

\gscore{ORC_SalveA}{}{Salut, Reine, mère de miséricorde, vie, douceur, espérance des hommes, salut ! Enfants d’Ève, nous crions vers toi dans notre exil. Vers toi, nous soupirons parmi les cris et les pleurs de cette vallée de larmes. Ô toi, notre avocate, tourne vers nous ton regard plein de bonté. Et à l’issue de cet exil, montre-nous Jésus, le fruit béni de tes entrailles. Ô clémente, ô bonne, ô douce Vierge Marie.}

\rubrica{ton simple:}
\gscore{ORC_SalveB}{}{}

\versiculus{Ora pro nobis, sancta Dei Génitrix.}{Ut digni efficiámur promissiónibus Christi.}{Prie pour nous, sainte Mère de Dieu.}{Afin que nous soyons rendus dignes des promesses du Christ.}
\versiculus{Orémus.\\
Omnípotens sempitérne Deus, qui gloriósæ Vírginis Matris Maríæ corpus et ánimam, ut dignum Fílii tui habitáculum éffici mererétur, Spíritu Sancto cooperánte, præpa\sylac{rá}sti:\pscross{} da, ut, cujus commemoratió\sylprep{ne} \sylprep{læ}\sylac{tá}mur,\psstar{} ejus pia intercessióne, ab instántibus malis et a morte perpétua liberémur. Per eúmdem Christum Dóminum nostrum.}{Amen.}{Prions.\\
Dieu tout-puissant et éternel, qui par la coopération du Saint-Esprit, as fait du corps et de l’âme de la glorieuse Vierge Marie  une demeure digne de ton Fils, accorde-nous, comme nous la célébrons dans la joie, d’être par son intercession, délivrés de tous les maux et de la mort éternelle. Par le même Jésus, le Christ, notre Seigneur.}{Amen.}

\versiculus{Divínum auxílium \cc máneat semper nobíscum.}{Amen.}{Que le secours divin demeure toujours avec nous.}{Amen.}

\end{document}